\section{Introduction}\label{sec:intro}

Every demand forecasting system trained on retail data faces the same quiet corruption: the training data records \emph{sales}, but the decisions it feeds, ordering, assortment, replenishment, need \emph{demand}. The two diverge exactly when it hurts. On a stockout day the till records zero however many customers wanted the item; on a day the shelf ran out at noon it records half a day's appetite; on a day demand exceeded the delivered quantity it records the delivery. A model fit to raw sales learns these artifacts as if they were demand patterns: it learns that the item "does not sell" after stockout runs and that peak demand "tops out" at shelf capacity. The forecast then feeds the ordering decision, the low forecast produces a low order, and the low order produces the next stockout. Yesterday's stockouts are baked into tomorrow's orders.

The repair is conceptually uniform even though its implementations differ. Treat the latent demand as the modeled quantity, treat whatever is recorded about availability as data, and make the corruption an explicit, switchable component of the joint model. The sales forecast is the model with the corruption component on; the demand forecast, the object planning actually needs, is the same model with the component switched off. What varies across situations is which corruption component the recorded data can support, and that is the organizing question of this paper: \emph{what do you observe about the stockout?}

This paper is a practitioner's guide built around three worked case studies, one per answer to that question, all implemented in the open-source \pkg{numpyro\_forecast} package and reproducible from its example notebooks (\S\ref{sec:repro}). When only a \emph{binary on-shelf flag} is recorded, an intermittent-demand smoother can gate its occurrence updates on availability, so off-shelf zeros stop eroding the demand estimate (\S\ref{sec:cs1}). When a \emph{per-observation censoring point} is recorded, a Tobit-style censored likelihood keeps capped days in the fit as lower bounds instead of corrupted values (\S\ref{sec:cs2}). When availability is \emph{fractional and noisily labeled}, a multiplicative saturating factor with a learned floor attaches to a hierarchical state-space model and survives label noise at a scale of $50{,}000$ series (\S\ref{sec:cs3}). Alongside the mechanisms, the case studies apply a common evaluation discipline (\S\ref{sec:scoring}): aggregate accuracy against recorded sales rewards models that reproduce the corruption, so demand models must be judged on calibration, stress-day behavior, and counterfactual plausibility instead, and one case study exhibits the trap concretely by showing the mechanism-free model winning every point metric against sales while halving its accuracy advantage into a deficit on the days demand exceeded supply.

The scope is deliberately bounded. This is not a benchmark: no mechanism is applied to another's data situation, so the paper characterizes mechanisms and their failure modes rather than ranking them, and the controlled comparison that a ranking would require is spelled out in \S\ref{sec:discussion}. None of the mechanisms is claimed as new; gating, censored likelihoods, and attenuation factors all have long histories reviewed in \S\ref{sec:related}, and the availability-gated smoother of \S\ref{sec:cs1} in particular is a one-gate modification of the method of \citet{teunter2011intermittent}. The contribution is the unified presentation, the honest accounting of what each mechanism does and does not buy on worked examples with stored, verifiable outputs, and the demonstration that the whole toolkit runs at production scale in a modern probabilistic-programming stack.
